\subsection{The category of Daniell closed integration systems}

\begin{definition}\label{def:daniell-closed}
  We say that an integration system $(X,L,I)$ on a set $X$ is
  \textbf{Daniell closed} if
  \[
    L^1=L .
  \]
\end{definition}

\begin{remark}[on terminology]\label{rem:why-daniell-closed}
  This condition is often called ``completeness''.
  We assign a different word to it because
  \textbf{``complete'' occurs in this paper in three senses}:
  \begin{itemize}
    \item $L^1=L$ (the present definition);
    \item completeness with respect to the norm $\|\cdot\|_I$
      (Theorem~\ref{thm:l1-norm-complete});
    \item completion of a measure (Subsection~\ref{sec:why-complete},
      Remark~\ref{rem:separation-not-sufficient}).
  \end{itemize}
  These are three distinct conditions, and in this paper all three live in the
  same subsection.
  We therefore call the first one \textbf{Daniell closed}.
  The name follows \cite{CWS} and \cite{CFW}, which use the word ``closure'' of a
  Daniell space.

  \textbf{It also clashes with ``complete'' in other books.}
  \cite{CWS} uses ``complete'' for the sandwich condition ($f\le h\le g$
  together with $\ell(f)=\ell(g)$ implies that $h$ belongs as well), which is
  none of the three above but rather our
  Definition~\ref{def:mod-closed}.
  That book further uses ``complete'' for order-complete lattices and for the
  completion of a measure space.
  The correspondence of terms is collected in
  Remark~\ref{rem:extended-real-collapse}.

  In \textbf{the title}, however, \textbf{we follow custom and say
  ``completeness''}.
  Outside this paper that is the standard word for $L^1=L$, and the coexistence
  of the three senses is an internal matter of this paper.
  In the body and the abstract we use ``Daniell closed'' consistently.
\end{remark}

\begin{remark}\label{rem:complete-implies-mct}
  A Daniell closed integration system automatically satisfies the monotone
  convergence theorem.
  Indeed, the Daniell extension $(X,L^1,I^1)$ satisfies the monotone convergence
  theorem by Lemma~\ref{lem:daniell-ext-satisfies-mct}, and since $L^1=L$ and
  $I^1=I$ (Lemma~\ref{lem:daniell-ext-contains-source}), the system $(X,L,I)$
  itself satisfies the monotone convergence theorem.

  We may therefore define $\mathbf{DInt}$ as a full subcategory of
  $\mathbf{MInt}$.
  Where we list the monotone convergence theorem and Daniell closedness side by
  side below, it is in order to make the relation to $\mathbf{MInt}$ easy to see.
\end{remark}

\begin{definition}
  The full subcategory of the category $\mathbf{MInt}$ of integration systems
  satisfying the monotone convergence theorem consisting of the Daniell closed
  integration systems is written
  \[
    \mathbf{DInt}
  \]
  and called the \textbf{category of Daniell closed integration systems}.
\end{definition}

\begin{definition}\label{def:inclusion-functors}
  We write the inclusion functors among the three categories above as
  \begin{align*}
    \iota&:\mathbf{MInt}\to\mathbf{Int},\\
    \iota_{D}&:\mathbf{DInt}\to\mathbf{Int},\\
    \iota_M&:\mathbf{DInt}\to\mathbf{MInt}
  \end{align*}
  respectively. All of them are inclusions of full subcategories, and
  $\iota_{D}=\iota\circ\iota_M$.
\end{definition}

\begin{lemma}\label{lem:daniell-ext-is-complete}
  Let $(X,L,I)$ be an integration system on a set $X$ and let
  $(X,L^1,I^1)$ be its Daniell extension.
  Then the integration system $(X,L^1,I^1)$ is Daniell closed.
\end{lemma}

\noindent
This statement itself is stated as a corollary in \cite{Yamagami} and is not new.
However, since this paper uses it essentially in the main theorem of
Section~\ref{chap:functoriality}, in the form ``the codomain of the
factorization is Daniell closed'', we give a proof in this subsection only.
\begin{proof}
  We have already checked that the monotone convergence theorem holds, so it
  suffices to show $(L^1)^1=L^1$.
  Since $L^1\subset(L^1)^1$ has already been shown
  (Lemma~\ref{lem:daniell-ext-contains-source}), it suffices to show
  $(L^1)^1\subset L^1$.

  Let $f\in (L^1)^1$. That is, there are sequences of functions
  $\{\varphi_n\},\{f_n\}\subset L^1$ such that
  \begin{enumerate}
    \item $|f_n|\le\varphi_n$;
    \item $\sum_{n=1}^\infty I^1(\varphi_n)<+\infty$;
    \item if $\sum_{n=1}^\infty \varphi_n(x)<+\infty$ then
      $f(x)=\sum_{n=1}^\infty f_n(x)$.
  \end{enumerate}
  For each $n$ put
  \[
    f_n^+:=f_n\lor0,
    \qquad
    f_n^-:=(-f_n)\lor0 .
  \]
  Since $L^1$ is a Riesz space we have $f_n^+,f_n^-\in L^1$, and
  \[
    f_n=f_n^+-f_n^-,
    \qquad
    |f_n|=f_n^++f_n^- .
  \]
  Moreover, by positivity of $I^1$ and $|f_n|\le\varphi_n$,
  \[
    \begin{aligned}
      \sum_{n=1}^\infty I^1(f_n^+)
      +\sum_{n=1}^\infty I^1(f_n^-)
      &=\sum_{n=1}^\infty I^1(|f_n|)\\
      &\le\sum_{n=1}^\infty I^1(\varphi_n)
      <+\infty.
    \end{aligned}
  \]

  Now define functions $p,q:X\to\mathbb{R}$ by
  \[
    p(x):=
    \begin{cases}
      \sum_n f_n^+(x) & \text{if $\sum_n f_n^+(x)$ converges},\\
      0 & \text{otherwise}
    \end{cases}
  \]
  and
  \[
    q(x):=
    \begin{cases}
      \sum_n f_n^-(x) & \text{if $\sum_n f_n^-(x)$ converges},\\
      0 & \text{otherwise.}
    \end{cases}
  \]
  By Lemma~\ref{lem:positive-series-flattening} we have $p,q\in L^1$.

  Put
  \[
    A:=\left\{x\in X\mathrel{}\middle|\mathrel{}\sum_{n=1}^\infty\varphi_n(x)<+\infty\right\} .
  \]
  If $x\in A$ then
  \[
    \sum_{n=1}^\infty|f_n(x)|
    \le
    \sum_{n=1}^\infty\varphi_n(x)
    <+\infty ,
  \]
  so the two series defining $p(x)$ and $q(x)$ both converge. Hence, by the
  condition on the original series representation,
  \[
    f(x)=\sum_{n=1}^\infty f_n(x)=p(x)-q(x)
    \qquad(x\in A).
  \]
  Putting
  \[
    r:=f-(p-q) ,
  \]
  the function $r$ vanishes on $A$.

  Choose $\varepsilon_n>0$ for each $n$ so that
  \[
    \sum_{n=1}^\infty\varepsilon_n<+\infty .
  \]
  Since $\varphi_n\ge0$, by Lemma~\ref{lem:positive-good-representation} we may
  choose for each $n$ a series representation
  \[
    \varphi_n\overset{\psi_{n,k}}{\simeq}\sum_k g_{n,k}
  \]
  of $\varphi_n$ satisfying
  \[
    \sum_{k=1}^\infty I(\psi_{n,k})
    \le I^1(\varphi_n)+\varepsilon_n .
  \]
  Hence
  \[
    \sum_{n,k}I(\psi_{n,k})
    \le
    \sum_n I^1(\varphi_n)+\sum_n\varepsilon_n
    <+\infty.
  \]

  For each $n$,
  \[
    0\overset{\psi_{n,k}}{\simeq}\sum_k0
  \]
  is a series representation of the function constantly equal to $0$. Put
  \[
    B:=
    \left\{
      x\in X
      \mathrel{}\middle|\mathrel{}
      \sum_{n,k}\psi_{n,k}(x)<+\infty
    \right\} .
  \]
  If $x\in B$ then $\sum_k\psi_{n,k}(x)<+\infty$ for each $n$, so
  \[
    \varphi_n(x)=\sum_k g_{n,k}(x)
  \]
  and
  \[
    0\le\varphi_n(x)
    \le\sum_k|g_{n,k}(x)|
    \le\sum_k\psi_{n,k}(x).
  \]
  Hence
  \[
    \sum_n\varphi_n(x)
    \le\sum_{n,k}\psi_{n,k}(x)
    <+\infty.
  \]
  Therefore $x\in A$, and
  \[
    r(x)=0=\sum_n0
    \qquad(x\in B).
  \]
  Applying Lemma~\ref{lem:flatten-series-representations} with $u_n=0$ and
  $u=r$, we get $r\in L^1$.

  Since $p,q,r\in L^1$ and $L^1$ is a vector space,
  \[
    f=p-q+r\in L^1 .
  \]
  Therefore $(L^1)^1\subset L^1$, and $(L^1)^1=L^1$ is proved.
\end{proof}

\subsubsection{Splitting Daniell closedness into two closure conditions}

Daniell closedness $L^1=L$ is a single closure condition passing through the
extension by series representations.
In this subsubsection we show that this condition splits into two independent
conditions, \textbf{closure under the monotone convergence theorem} and
\textbf{closure under modification on null sets}
(Theorem~\ref{thm:complete-intrinsic}).
That neither condition can be dropped is seen in
Remark~\ref{rem:two-conditions-independent}.

The two conditions on the right are conditions that occur independently in
measure theory.
The monotone convergence theorem corresponds directly, and closure under
modification on null sets corresponds to completeness of the measure.
This decomposition is important for gauging the content of the main theorem of
Section~\ref{chap:functoriality}
(Remark~\ref{rem:main-theorem-not-tautology}).

\begin{definition}\label{def:null-set}
  Let $(X,L,I)$ be an integration system on a set $X$.
  We say that a subset $Z\subset X$ is an \textbf{$I$-null set} if there exists a
  sequence of functions $\{g_n\}$ in $L$ satisfying the following three
  conditions.
  \begin{enumerate}
    \item $g_n\ge0$ for every $n$;
    \item $\sum_{n=1}^{\infty}I(g_n)<+\infty$;
    \item $\sum_{n=1}^{\infty}g_n(x)=+\infty$ for every $x\in Z$.
  \end{enumerate}
\end{definition}

\noindent
This definition does not mention $L^1$.
Note that it is written using only the data given by $(X,L,I)$.
The set of all points at which condition (3) of a series representation
(Definition~\ref{def:series-representation}) does not constrain the value is
always an $I$-null set, as one sees by taking $\{g_n\}$ to be the dominating
functions $\{\varphi_n\}$.

\begin{lemma}\label{lem:null-sets-countable-union}
  A countable union of $I$-null sets is an $I$-null set.
\end{lemma}
\begin{proof}
  Let $Z_1,Z_2,\dots$ be $I$-null sets, choose for each $Z_k$ a sequence
  $\{g^{(k)}_n\}_n$ as in Definition~\ref{def:null-set}, and put
  $C_k:=\sum_nI(g^{(k)}_n)<+\infty$.
  Put $a_k:=2^{-k}/(1+C_k)>0$ and $\tilde{g}^{(k)}_n:=a_kg^{(k)}_n\in L$.
  Arranging the family $\{\tilde{g}^{(k)}_n\}_{k,n}$ into a single sequence,
  every term is nonnegative and
  \[
    \sum_{k,n}I(\tilde{g}^{(k)}_n)=\sum_ka_kC_k\le\sum_k2^{-k}=1<+\infty .
  \]
  Moreover, if $x\in\bigcup_kZ_k$ then $x\in Z_k$ for some $k$ and
  \[
    \sum_n\tilde{g}^{(k)}_n(x)=a_k\sum_ng^{(k)}_n(x)=+\infty .
  \]
  The remaining terms are nonnegative, so the sum of the rearranged sequence is
  $+\infty$ as well.
\end{proof}

\begin{lemma}\label{lem:vanish-off-null-set}
  Let $(X,L,I)$ be an integration system and let $Z\subset X$ be an $I$-null set.
  If $g\in L$ vanishes on $X\setminus Z$ then
  \[
    I(|g|)=0,
    \qquad\text{and in particular}\qquad
    I(g)=0 .
  \]
\end{lemma}
\begin{proof}
  Choose $\{g_n\}$ as in Definition~\ref{def:null-set} and put
  $C:=\sum_nI(g_n)<+\infty$.
  Take an arbitrary $\varepsilon>0$ and put
  \[
    u_N:=|g|\land\Bigl(\varepsilon\sum_{n\le N}g_n\Bigr) .
  \]
  Since $L$ is a Riesz space we have $u_N\in L$.

  We check that $u_N$ increases monotonically to $|g|$ pointwise.
  If $x\notin Z$ then $|g|(x)=0$, and the second term is nonnegative, so
  $u_N(x)=0=|g|(x)$.
  If $x\in Z$ then $\varepsilon\sum_{n\le N}g_n(x)\uparrow+\infty$, so
  $u_N(x)=|g|(x)$ for $N$ large enough, while before that
  $u_N(x)=\varepsilon\sum_{n\le N}g_n(x)$, which is nondecreasing.
  In either case $u_N(x)\uparrow|g|(x)$.

  Since $|g|\in L$, Lemma~\ref{lem:integration-is-continuous} gives
  $I(u_N)\uparrow I(|g|)$.
  On the other hand $u_N\le\varepsilon\sum_{n\le N}g_n$, so
  Lemma~\ref{lem:integration-preserve-order} gives
  \[
    I(u_N)\le\varepsilon\sum_{n\le N}I(g_n)\le\varepsilon C.
  \]
  Hence $I(|g|)\le\varepsilon C$. As $\varepsilon>0$ was arbitrary, $I(|g|)=0$,
  and $|I(g)|\le I(|g|)=0$ by
  Lemma~\ref{lem:triangle-ineq-for-integration}.
\end{proof}

\begin{definition}\label{def:mod-closed}
  We say that an integration system $(X,L,I)$ on a set $X$ is
  \textbf{closed under modification on null sets} if
  \[
    h|_{X\setminus Z}=f|_{X\setminus Z}
    \quad\Longrightarrow\quad
    h\in L
  \]
  holds for every $f\in L$, every $I$-null set $Z\subset X$ and every function
  $h:X\to\mathbb{R}$.
\end{definition}

\noindent
On $Z$ the function $h$ is subject to no constraint at all; only that it is
real-valued.

\begin{theorem}[decomposition of Daniell closedness]\label{thm:complete-intrinsic}
  Let $(X,L,I)$ be an integration system on a set $X$. The following are
  equivalent.
  \begin{enumerate}
    \item[(i)] $(X,L,I)$ is Daniell closed, that is, $L^1=L$.
    \item[(ii)] $(X,L,I)$ satisfies the monotone convergence theorem and is
      closed under modification on null sets.
  \end{enumerate}
\end{theorem}
\begin{proof}
  \textbf{(i)$\Rightarrow$(ii)}.
  The monotone convergence theorem was seen in
  Remark~\ref{rem:complete-implies-mct}.
  We show closure under modification on null sets.
  Let $f\in L$, let $Z$ be an $I$-null set, and let $h:X\to\mathbb{R}$ agree with
  $f$ on $X\setminus Z$.
  Choose $\{g_n\}$ as in Definition~\ref{def:null-set} and put
  \[
    f_1:=f,\quad f_n:=0\ (n\ge2),
    \qquad
    \varphi_1:=|f|,\quad \varphi_{n+1}:=g_n\ (n\ge1) .
  \]
  Then $|f_n|\le\varphi_n$ is clear, and
  \[
    \sum_{n}I(\varphi_n)=I(|f|)+\sum_nI(g_n)<+\infty .
  \]
  Moreover, if $\sum_n\varphi_n(x)<+\infty$ then $\sum_ng_n(x)<+\infty$, so
  $x\notin Z$ and therefore $h(x)=f(x)=\sum_nf_n(x)$.
  Hence $\{f_n\},\{\varphi_n\}$ is a series representation of $h$, and
  $h\in L^1=L$.

  \textbf{(ii)$\Rightarrow$(i)}.
  Since $L\subset L^1$ by
  Lemma~\ref{lem:daniell-ext-contains-source}, it suffices to show
  $L^1\subset L$.

  Let $f\in L^1$ and choose a series representation
  $f\overset{\varphi_n}{\simeq}\sum f_n$.
  Put
  \[
    Z:=\Bigl\{x\in X\;\Bigm|\;\sum_n\varphi_n(x)=+\infty\Bigr\} .
  \]
  From $|f_n|\le\varphi_n$ we get $\varphi_n\ge0$, and
  $\sum_nI(\varphi_n)<+\infty$, so $\{\varphi_n\}$ witnesses that $Z$ is an
  $I$-null set.

  Put $p_n:=f_n\lor0$ and $m_n:=(-f_n)\lor0$, so that $p_n,m_n\in L$ and
  \[
    f_n=p_n-m_n,
    \qquad
    0\le p_n,m_n\le|f_n|\le\varphi_n .
  \]
  Put
  \[
    P_N:=\sum_{n\le N}p_n\in L,
    \qquad
    M_N:=\sum_{n\le N}m_n\in L .
  \]
  Both are nondecreasing in $N$, and
  Lemma~\ref{lem:integration-preserve-order} gives
  \[
    I(P_N)\le\sum_{n\le N}I(\varphi_n)\le\sum_nI(\varphi_n)<+\infty ,
  \]
  and likewise for $M_N$.

  If $x\notin Z$ then $\sum_n\varphi_n(x)<+\infty$, so $\sum_np_n(x)$ and
  $\sum_nm_n(x)$ both converge. We therefore define
  \[
    P'(x):=
    \begin{cases}
      \sum_np_n(x) & (x\notin Z)\\
      0 & (x\in Z)
    \end{cases}
  \]
  and define $M'$ similarly. Both are real-valued functions on $X$.
  Let $P_N'$ be the function agreeing with $P_N$ on $X\setminus Z$ and vanishing
  on $Z$, and define $M_N'$ similarly.
  By hypothesis $(X,L,I)$ is closed under modification on null sets, so
  $P_N',M_N'\in L$.

  Since $P_N-P_N'$ vanishes on $X\setminus Z$,
  Lemma~\ref{lem:vanish-off-null-set} gives $I(P_N-P_N')=0$, that is,
  \[
    I(P_N')=I(P_N)\le\sum_nI(\varphi_n)<+\infty .
  \]
  Moreover $P_N'(x)\uparrow P'(x)$ for every $x\in X$ (for $x\notin Z$ by
  definition, and for $x\in Z$ both sides are $0$).
  By hypothesis $(X,L,I)$ satisfies the monotone convergence theorem, so
  $P'\in L$. Similarly $M'\in L$.

  Hence $P'-M'\in L$.
  If $x\notin Z$ then $\sum_nf_n(x)$ converges absolutely by
  Lemma~\ref{lem:series-exp-is-abs-conv}, so
  \[
    (P'-M')(x)=\sum_np_n(x)-\sum_nm_n(x)=\sum_nf_n(x)=f(x) .
  \]
  That is, $f$ agrees with the element $P'-M'$ of $L$ on $X\setminus Z$.
  Using closure under modification on null sets again, $f\in L$.
\end{proof}

\medskip
\noindent
\textbf{On the construction in (ii)$\Rightarrow$(i).}
The functions $P'$ and $M'$ built in the proof above are both monotonically
nondecreasing limits of elements of $L$ with bounded integrals, that is,
\textbf{upper functions} in the classical sense.
What the proof exhibits is therefore the representation of an $f\in L^1$ as
agreeing, off a null set, with the difference $P'-M'$ of two upper functions.
Constructing the Lebesgue integral as differences of upper functions is a
classical method, and a representation of this shape is well known in that
context.
We do not use this correspondence below, and so cite no particular source for
it.

\begin{corollary}\label{cor:mintbar-intrinsic}
  $\mathbf{DInt}$ is the full subcategory of $\mathbf{Int}$ consisting of the
  \textbf{integration systems that satisfy the monotone convergence theorem and
  are closed under modification on null sets}.
\end{corollary}
\begin{proof}
  By Theorem~\ref{thm:complete-intrinsic}.
\end{proof}

\begin{remark}\label{rem:not-about-mentioning-L1}
  The description above does not pass through the construction of the Daniell
  extension, but \textbf{that in itself is not where the point lies}.
  The condition $L^1=L$, when unfolded, also reads
  \begin{quote}
    every function admitting a series representation by elements of $L$ belongs
    to $L$,
  \end{quote}
  which uses nothing but $L$, $I$ and pointwise conditions.
  ``Whether $L^1$ is mentioned'' does not distinguish the two.

  What matters is that one unfamiliar closure condition splits into two
  conditions that occur independently in measure theory.
  We return to this point in
  Remark~\ref{rem:main-theorem-not-tautology}.
\end{remark}

\begin{example}
  Let us check this in the example of
  Remark~\ref{rem:D-I-not-left-adjoint} ($X=\{0,1,2\}$, $L=\{(a,b,b)\}$,
  $I(a,b,b)=a$).
  The function $g_n:=(0,1,1)\in L$ satisfies $g_n\ge0$ and $I(g_n)=0$, and
  $\sum_ng_n(x)=+\infty$ at $x=1,2$, so $Z=\{1,2\}$ is an $I$-null set.
  But $h=(0,1,0)$, obtained from $f=(0,0,0)\in L$ by modifying it on $Z$, does
  not belong to $L$.
  Hence this integration system is not closed under modification on null sets
  and, by Theorem~\ref{thm:complete-intrinsic}, is not Daniell closed.
  This agrees with the direct computation $L^1=\mathbb{R}^X\ne L$ in that
  example.

  On the other hand, in Example~\ref{ex:stone-fails} ($X=\{1,2\}$,
  $L=\{(a,2a)\}$, $I(a,2a)=a$), if $g_n=(c_n,2c_n)\ge0$ satisfies
  $\sum_nI(g_n)=\sum_nc_n<+\infty$ then $\sum_ng_n(x)<+\infty$ holds at both
  $x=1$ and $x=2$.
  Hence the only $I$-null set is the empty set, and the condition on
  modification on null sets holds vacuously.
  The monotone convergence theorem holds as well, so this integration system is
  Daniell closed by Theorem~\ref{thm:complete-intrinsic}, in agreement with
  $L^1=L$ as stated in that example.
\end{example}

\begin{remark}[neither condition can be dropped]\label{rem:two-conditions-independent}
  The two conditions in (ii) of Theorem~\ref{thm:complete-intrinsic} are
  independent.

  \textbf{That the monotone convergence theorem cannot be dropped} is shown by
  Example~\ref{ex:counting} ($X=\mathbb{N}$, the functions vanishing outside
  finitely many $n$, and summation).
  In this integration system, $g\in L$ with $g\ge0$ satisfies $g(x)\le I(g)$ for
  every $x$, so for a nonnegative sequence $\{g_n\}\subset L$ with
  $\sum_nI(g_n)<+\infty$ we have $\sum_ng_n(x)<+\infty$ at every $x$.
  Hence the only $I$-null set is the empty set, and the condition on
  modification on null sets holds vacuously.
  On the other hand
  \[
    f_N:=\sum_{k\le N}2^{-k}1_{\{k\}}\in L
  \]
  satisfies $I(f_N)\le1$ while increasing monotonically pointwise to
  $f(k)=2^{-k}$, and $f\notin L$ because $f$ does not have finite support.
  That is, the monotone convergence theorem fails.
  Indeed this integration system is not Daniell closed, and $L^1=\ell^1$.

  \textbf{That modification on null sets cannot be dropped} is shown by the
  example of Remark~\ref{rem:D-I-not-left-adjoint}.
  The integration system of that example satisfies the monotone convergence
  theorem but, as seen in the example above, is not closed under modification on
  null sets, hence not Daniell closed.
\end{remark}

\begin{remark}
  The two conditions in (ii) of Theorem~\ref{thm:complete-intrinsic} correspond
  to two properties that $L^1(\mu)$ satisfies in measure theory.
  The monotone convergence theorem corresponds directly, and closure under
  modification on null sets corresponds to the fact that $L^1(\mu)$ contains the
  functions measurable with respect to the completion of $\mu$.
  In Proposition~\ref{prop:l1-decomposition} of
  Subsection~\ref{sec:why-complete} we will see that the gap between the two
  reflections amounts exactly to the functions supported on null sets; the
  present theorem explains the reason for that from the side of Daniell
  closedness.
  That is, \textbf{to be an object of $\mathbf{DInt}$ one must be closed under
  modification on null sets}.
\end{remark}

\begin{remark}[the decomposition collapses in the extended-real formalism]\label{rem:extended-real-collapse}
  That the two conditions of Theorem~\ref{thm:complete-intrinsic} are
  independent \textbf{depends on this paper's restriction of functions to real
  values}.
  If the range is widened to
  $\overline{\mathbb{R}}=[-\infty,+\infty]$, then condition (b) follows from
  condition (a).

  Indeed \cite{CWS} adopts that formalism.
  It defines the counterpart of a Riesz space inside
  $\overline{\mathbb{R}}^X$ under the name Riesz lattice, and calls the triple
  corresponding to our integration system a Daniell space.
  It then gives the two conditions separate names.
  \begin{itemize}
    \item Its Definition 1.4.1 calls a triple \textbf{closed} if, for every
      monotone sequence whose integrals are bounded, the pointwise limit again
      belongs to the space and the integral commutes with the limit.
      This is our Definition~\ref{def:satisfies-mct} (satisfying the monotone
      convergence theorem).
    \item Its Definition 1.4.7 calls a triple \textbf{complete} if
      $f\le h\le g$ together with $\ell(f)=\ell(g)$ implies that $h$ belongs to
      the space as well.
      This sandwich condition is in substance our
      Definition~\ref{def:mod-closed}.
  \end{itemize}
  Its Corollary 1.4.8, however, states that
  \begin{quote}
    every closed Daniell space is complete.
  \end{quote}
  That is, \textbf{in the extended-real formalism (a)$\Rightarrow$(b), and the
  two conditions are not independent.}

  The reason is the following. Let $Z$ be an $I$-null set and take the sequence
  $\{g_n\}$ of Definition~\ref{def:null-set}.
  The partial sums $\sum_{n\le N}g_n$ are monotonically nondecreasing with
  bounded integrals, so under (a) their pointwise limit belongs to the space.
  If extended-real values are allowed, that limit is an element of the space as
  a \emph{function} taking the value $+\infty$ on $Z$.
  Null sets can then be named from inside the space, as the sets where an
  element diverges to $\infty$, and (b) follows from the axioms of a Riesz
  lattice, which originate in the fact that $\infty-\infty$ is undefined.

  \textbf{This argument does not go through here.}
  Since $L\subset\mathbb{R}^X$, the limit above cannot be an element of $L$.
  A null set cannot be captured from inside $L$ as divergence to $\infty$, and
  it is precisely for this reason that
  Definition~\ref{def:null-set} takes the form of a series and that condition
  (b) remains as an independent condition.
  The counterexample of
  Remark~\ref{rem:two-conditions-independent} ($X=\{0,1,2\}$,
  $L=\{(a,b,b)\}$, $I(a,b,b)=a$) was an integration system satisfying the
  monotone convergence theorem but not closed under modification on null sets.
  That is, it is also a \textbf{counterexample to the real-valued form of
  Corollary 1.4.8}.

  This is not a question of which formalism is better.
  The correspondence is that \textbf{widening the range makes null sets nameable
  inside the space, at the cost of collapsing the two closure conditions into
  one}.
  Theorem~\ref{thm:complete-intrinsic} is a decomposition visible only in a
  formalism restricted to real values.
\end{remark}
